% !TeX root = ../master.tex

\chapter{Methodology}
\label{ch:methodology}

% Shared experimental setup across all three parts

\section{Preprocessing: Context Extraction}
\label{sec:preprocessing}

% - Argument definition extraction: PDF -> GPT with structured output (Pydantic schemas)
%   - 3--10 sentence definition per dataset, grounded in dataset's own theoretical framing
%   - Excludes logistics, collection methodology, experimental results
% - Annotation guideline extraction: 4 datasets (ABSTRCT, ARGUMINSCI, PE, USELEC)
%   - Synthesize decision rules an annotator would apply
%   - Remaining 6 datasets: "Not available" placeholder
% - Document context extraction: 6 datasets with source documents
%   - 2 preceding sentences via punctuation-based splitting + string matching
%   - Per-sample files for loading at inference time

\section{Models}
\label{sec:models}

% - Five decoder-based LLMs spanning 7B to frontier scale:
%   - Mistral-7B-Instruct-v0.2 (7B) -- small open-weight baseline
%   - Llama-3.1-8B-Instruct (8B) -- comparable size, different architecture family
%   - Mistral-Small-24B-Instruct (24B) -- mid-range
%   - Llama-3.1-70B-Instruct (70B) -- large open-weight
%   - GPT-4.1 (frontier-scale, closed-weight) -- ceiling reference
% - All accessed via OpenAI-compatible API (multi-provider: Portkey/Azure, Together AI, Groq, Ollama)
% - Temperature=0 for reproducibility, max_tokens constrained

\section{Prompt Design}
\label{sec:prompt_design}

% - Two-message structure: system prompt + user prompt
% - System prompt: role (dataset annotator), task (classify as Argument/No-Argument),
%   constraint (respond with exactly one label), rule (only "Argument" if clearly matches definition)
% - Context injected into system prompt (definition, guideline, document context)
% - User prompt: only the target sentence
% - Same template across all models, conditions, datasets -- only context block varies
% - Output normalization: map LLM output to exactly "Argument" or "No-Argument"

\section{Text Manipulations}
\label{sec:manipulations}

% - M0 (Original): unmodified sentence
% - M1 (Feger): remove stop words, function words, discourse markers, punctuation
%   - Implemented via spaCy, following Feger et al. exactly
%   - Removes ~50% of words, preserving only content words
% - M2 (Shuffle): randomly permute word order (seed=42)
%   - All words preserved, only order destroyed
%   - Complementary diagnostic: if order matters, model uses structural cues
% - Rationale: M1 tests reliance on non-content words, M2 tests reliance on word order

\section{Evaluation Metrics}
\label{sec:metrics}

% - Primary: Macro F1 per dataset and model
% - Delta_feger = F1(M0) - F1(M1): sensitivity to content-word-only input
% - Delta_shuffle = F1(M0) - F1(M2): sensitivity to word order destruction
% - Comparison baseline: encoder delta <= 0.02 from Feger et al.
% - Balanced sampling: n/2 per class, deterministic ordering
% - Per-dataset and overall (mean across datasets) reporting

\section{Contamination Control}
\label{sec:contamination_control}

% - Sentence-completion test following Golchin & Surdeanu (ICLR 2024)
% - Split sentences at midpoint, prompt model to complete
% - Metrics: ROUGE-L and exact match rates
% - Threshold: ROUGE-L > 0.5 flags contamination
% - Temporal analysis: older (pre-2018) vs. newer datasets
% - Results from contamination_test_outputs inform interpretation of all experiments
